% ----- Consignes exo 5 ----- %
\begin{td-exo}[Autour du problème d'ordonnancement avec contraintes de précédence]\,\\ % 5 
	On rappelle que le problème \(P \ \mid \text{prec}\ \mid C_{\max}\) est défini de la manière suivante:

	\begin{problem}{Le problème d'ordonnancement \(P \ \mid\ \text{prec}\ \mid\ C_{\max}\)}
    \Input{\(n\) tâches de durées respectives \(p_1, p_2, \ldots, p_n\) et \(\impliedby (I, U)\) une relation résumant les contraintes de précédence entre les tâches.}
    \Question{Ordonnancer ces tâches sur \(m\) machines identiques en une durée totale minimale notée \(C_{\max}\) tout en respectant les contraintes de précédence.}
\end{problem}

	\begin{enumerate}
		\item On veut montrer que pour n'importe quel ordonnancement de liste \(LS\) on a
		\begin{equation*}
			\frac{C_{\max}^{LS}}{C_{\max}^\ast} \leq 2 - \frac{1}{m}
		\end{equation*}

		\begin{enumerate}
			\item Pour cela, considérons la tache \(T_{i_1}\) dont la fin d'exécution coincide avec \(C_{\max}^{LS}\). 
			Soit \(T_{i_2} = T_j \setminus T_j \in \Gamma^-(T_{i_1}) \land C_j = \max_{T_j \in \Gamma^-(T_{i_1})} C_i\) où \(C_i = t_i + p_i\) désigne la fin d'exécution de la tache \(i\) et \(\Gamma^-(T_{i_1})\) désigne l'ensemble des prédécesseurs de \(i_j\).

			En procédant ainsi de suite, on a \(T_{i_k} = T_j \setminus T_j \in \Gamma^-(T_{i_{k-1}}) \land C_j = \max_{T_j \in \Gamma^-(T_{i_{k-1}})} CT_i\). On appellera \(P\) le chemin constitué de ces taches \(T_{i_1}, T_{i_2}, \ldots, T_{i_k}\).

			Est-ce qu'il existe des processeurs (ou machines) inactifs entre la fin d'exécution d'un prédécesseur de \(T_{i_1}\) et le début d'exécution de \(T_{i_1}\)?

			Est-ce qu'il existe des processeurs (ou machines) inactifs entre la fin d'exécution d'un prédécesseur de \(T_{i_k}\) et le début d'exécution de \(T_{i_k}, \forall k\)?
			
			Que pouvons-nous conclure sur les dates de début d'exécution des taches appartant au chemin \(P\)?

			\item A partir de l'équation suivante, conclure sur le ratio:
			\begin{equation*}
				C_{\max}^{LS} = \frac{T_{\text{seq}} + T_{\text{idle}}}{m}
			\end{equation*}

			\item Pouvons-nous appliquer le même algorithme au problème \(P \ \mid \text{prec}, r_j \mid C_{\max}\), c'est-à-dire le même problème décrit ci-dessus mais avec le début des taches \(T_j\) qui ne peut se faire qu'après la date de disponibilité \(r_j\)?
		\end{enumerate}

		\item Montrer que la borne supérieure est atteinte.
		Pour cela, considérer l'instance suivante:
		\begin{itemize}
			\item \(m\) taches indépendantes,
			\item la tache \(m\) précède les taches \(m+1, m+2, \ldots, 2m\),
			\item la tache \(2m + 1\) suit les taches \(m+1, m+2, \ldots, 2m\).
		\end{itemize}
		Les durées des taches sont les suivantes:
		\begin{itemize}
			\item \(p_1 = p_2 = \cdots = p_{m-1} = m - 1\),
			\item \(p_m = \varepsilon\),
			\item \(p_{m+1} = p_{m+2} = \cdots = p_{2m} = 1\),
			\item \(p_{2m + 1} = m - 1\).
		\end{itemize}
	\end{enumerate}
\end{td-exo}

% ----- Solutions exo 5 ----- %
\iftoggle{showsolutions}{ 
	\begin{td-sol}[]\ % 5
		\begin{enumerate}
			\item Sur le ratio d'ordonnancement de liste:
			\begin{enumerate}
				\item On obtient un chemin incompressible, les dates d'exécution ne peuvent être avancées.

				\item On a bien 
				\begin{equation*}
					m\cdot C_{\max} = T_{\text{seq}} + T_{\text{idle}}
				\end{equation*}
				d'où
				\begin{equation*}
					\begin{aligned}
						C_{\max} 
						&= \frac{T_{\text{seq}} + T_{\text{idle}}}{m} \\
						& \leq \frac{T_{\text{seq}} + (m-1) \sum_{i=1}^n P_i p}{m} \\
						& = OPT + \left( 1 - \frac{1}{m} \right) OPT \\
						& = \left( 2 - \frac{1}{m} \right) OPT
					\end{aligned}
				\end{equation*}
			\end{enumerate}
		\end{enumerate}
	\end{td-sol}
}{}
